\section{Injection resistance by matcher-delimiting untrusted input}
\label{sec:resist}

Delimiting untrusted input with a matcher pair rather than a nonmatcher (such as
a quote) prevents that input from breaking out of its slot, since escaping the
delimiter would require an unmatched matcher, which the always-embeddability
theorem of matchertext forbids in valid matchertext~\cite{matchertext}.
We develop this first for SQL injection, then for cross-site scripting, whose
several distinct HTML escaping contexts collapse to one discipline.
The same argument then covers two settings where the host already delimits with
matchers, so the discipline asks for no new syntax: LDAP directory filters,
whose parenthesised assertions are matcher-delimited already, and the PDF
document format, whose literal strings already require matcher balance.

\subsection{Preventing SQL injection}
\label{sec:resist:sql}

The discipline is worth making concrete in the setting that most sharply
motivates it, untrusted input embedded into a query language such as SQL.
Recall the classic vulnerability of \cref{sec:intro}: a server composes a clause
like \verb|"WHERE name = '"+userName+"'"|, and a \verb|userName| containing a
quote closes the literal early and appends attacker-controlled
SQL~\cite{clarke12sql}, the quote being dangerous precisely because it is a
\emph{nonmatcher} the discipline leaves unconstrained.
We instead delimit each untrusted value with the matcher pair
\verb|[|$\dots$\verb|]| and have the query processor treat the content as
embedded matchertext, ending it by matcher balance rather than by scanning for a
closing quote, so the clause becomes \verb|WHERE name = [|\textit{userName}\verb|]|.

Concretely, the canonical injection \verb|x' OR '1'='1| now fails two ways.
Supplied verbatim it contains no matchers, so it embeds as the literal name
\verb|x' OR '1'='1|, its quotes inert because nothing inside \verb|[|$\dots$\verb|]|
is read as SQL.
When adapted to attack the new delimiters, as in \verb|] OR 1=1 --|, it carries
an unmatched \verb|]| which is not matchertext, and thus cannot embed at all.
Because the guarantee depends only on the delimiter being a matcher, not on what
the value denotes, it holds wherever untrusted input appears, not only string
values but numbers and even identifiers such as table or column names.
That last case is a strength over parameterized queries (prepared statements),
whose placeholders bind only \emph{values}: an identifier cannot be parameterized
and is usually left to ad-hoc allow-listing, whereas \verb|FROM [|\textit{table}\verb|]|
covers it directly.
Which tables a caller may name remains an authorization question, distinct from
injection.

Free-form fields, a name, a comment, a search box, may legitimately contain
unmatched matchers (\eg a fragment \verb|f(x| or \verb|smile :]|), so rather than
reject them we pass them through a total encoder \textsc{ToMatchertext}$(v)$ that
escapes only the \emph{unmatched} matchers in $v$~\cite{matchertext}, and decode
on read.
Defined on every input, it never rejects, and because its escapes are themselves
balanced its output is always valid matchertext, its sole and host-independent
obligation, checkable in isolation by the \textsc{Verify} scan and saying nothing
about SQL.
The failure mode is therefore benign: an encoder bug yields mis-encoded
\emph{data}, never a breakout, where in conventional escaping one missed case
\emph{is} the injection.

Under a matchertext-aware host that ends each value by matcher balance,
matchertext offers a \emph{structural} property: no valid-matchertext value, in
any position, can escape its delimiters or alter the surrounding query structure,
closing off the sub-class of injection that depends on breaking out of a
delimiter.
This is a syntactic-closure property; it does not by itself establish that later
decoding, reparsing, or execution of the contained value is safe.
It does not adjudicate whether a well-formed value is \emph{authorized}, which
remains a matter of access control, and it complements rather than replaces
parameterized queries where those already suffice.

\subsection{Cross-site scripting across HTML contexts}
\label{sec:resist:xss}

The same discipline addresses cross-site scripting across three different
injection sites, each with its own escaping regime.
A value placed as element content, between \verb|<div>| and \verb|</div>|, breaks
out on \verb|<| or \verb|&|; the same value in a quoted attribute,
\verb|<input value="|\emph{v}\verb|">|, breaks out on the active quote rather than
\verb|<|, and in an unquoted attribute the breakout set widens to whitespace and
\verb|=|; a value in a script body, \verb|<script>var n="|\emph{v}\verb|"</script>|,
has two independent exits, the JavaScript string quote and the raw
\verb|</script>| sequence that ends the element regardless of JavaScript syntax.
The recurring bug is applying one escaper, usually HTML-entity encoding, in all
three contexts, where it is correct only for element content.

Matcher-delimiting collapses the three to a single discipline, exactly as for
SQL: each hole is read verbatim up to a matched close bracket and the host
performs no transformation, using matchertext markup forms~\cite{matchertext}
(\eg \verb|<div [|\emph{v}\verb|]>|, \verb|value=[|\emph{v}\verb|]|, and
\verb|<![MDATA[|\emph{v}\verb|]]>|).
Because the matchers are \verb|()[]{}| and not \verb|<>|, the characters that
drove every breakout above, \verb|<|, \verb|&|, the quote, and even the literal
\verb|</script>| sequence, become ordinary data inside a matched hole: the parser
counts brackets rather than scanning for tags, so an injected
\verb|<img src=x onerror=steal()>| is inert text.
The three context-specific sanitizers thus reduce to the single \textsc{Verify}
recognizer, the same one used for SQL, whose only rejection is an unmatched
matcher.

As with the numeric-context caveat for SQL, the guarantee is structural and stops
at sinks that execute their contents by design.
A payload kept dutifully inside \verb|<a href=[|\emph{v}\verb|]>| stays within the
attribute yet still runs when \emph{v} is a \verb|javascript:| URL, because that
is an execution context, not a structural breakout; an \verb|onerror| handler is
the same.
Such semantic sinks are out of scope.
A second limit is expressiveness: matchertext content is verbatim text and cannot
carry child markup, which is exactly right for embedding an untrusted string as
inert data but is not a substitute for rendering a user's rich HTML, a harder
problem left to allow-list sanitizers.
The clean guarantee covers the data-not-markup case, where the large majority of
injection bugs arise.

\subsection{Preventing LDAP injection}
\label{sec:resist:ldap}

LDAP directory queries are the sharpest illustration of the discipline, because
the language already delimits its structure with a matcher.
A search filter nests parenthesised assertions, \eg the bind check
\verb|(&(uid=|\textit{userName}\verb|)(userPassword=|\textit{pass}\verb|))|, and
the classic vulnerability~\cite{alonso08ldap} is a \verb|userName| that closes its
assertion early and injects further clauses, as in \verb|*)(uid=*|, so that the
password assertion no longer constrains the match and the bind succeeds against an
arbitrary account.
Where SQL forced us to \emph{introduce} a matcher pair, LDAP needs none: the
assertion's own parentheses \verb|(uid=|$\dots$\verb|)| already delimit the value,
so we simply have the directory processor end that value by matcher balance rather
than at the first \verb|)|.

The canonical injection then fails exactly as before: adapted to the delimiter,
\verb|*)(uid=*| carries an unmatched \verb|)|, is not matchertext, and cannot
embed at all.
A further strength is specific to LDAP: the discipline coincides with the
language's native syntax, and RFC~4515 value-encoding already escapes unbalanced
parentheses~\cite{rfc4515}, so a conformant emitter is already close to compliant
and needs no new hosting form.

\subsection{Preventing PDF object injection}
\label{sec:resist:pdf}

The same structural argument extends to the PDF document format, whose literal
strings are delimited by parentheses and, like the PostScript strings they descend
from, admit balanced parentheses unescaped, requiring an escape only for an
unmatched parenthesis or backslash~\cite{iso32000}: the format already mandates
matcher balance for its strings.
A server-side generator that concatenates untrusted input into a
\verb|(|$\dots$\verb|)| string, a \verb|/URI (|$\dots$\verb|)|, or a
\verb|/JS (|$\dots$\verb|)| action lets an attacker close the string and append
new objects or dictionary keys, adding annotations, actions, or JavaScript that
Acrobat and browser viewers will parse and run, the ``XSS for PDFs'' class that
yields script execution, request forgery, and document exfiltration~\cite{heyes20pdf}.

Emitting each untrusted value through \textsc{ToMatchertext} closes this by
construction.
Placed in \verb|/URI (|\textit{v}\verb|)|, the canonical breakout
\begin{center}\footnotesize
\verb|)>>/AA<</O<</S/JavaScript/JS(app.alert(1))>>>>(|
\end{center}
begins with an unmatched \verb|)|, is not matchertext, and is escaped to inert
data rather than embedded, while a balanced value such as \verb|f(x)| embeds
unchanged.
As with SQL and XSS the guarantee is structural and stops at sinks that execute by
design: data an application deliberately routes into a \verb|/JS| action still
runs, though it can no longer escape its string to \emph{introduce} one, which is
where the reported attacks overwhelmingly live.

\subsection{Fitting matchertext into existing languages}
\label{sec:resist:fitting}

\subsection{Comparison with existing defenses}
\label{sec:resist:baselines}

